% GENERATED FILE. Edit canonical lessons, then run npm run generate:books.
% canonical-source-hash: fnv1a64:686a9bfe2941d904
% canonical-lessons: LA-C260-exprimo, LA-C260-alo, LA-C260-compono, LA-C260-animadverto, LA-C260-erubesco

\chapter{To press out, To feed, To put together, To notice, To blush}
\label{ch:la-to-press-out-to-feed-to-put-together-to-notice-to-blush}
\hlchaptermodality{260}

\begin{chapteropening}
\textbf{By the end of this chapter:} \emph{I can say to press out, to feed, to put together, to notice and to blush.}

Complete the last lesson of chapter 260: I can say to press out, to feed, to put together, to notice and to blush.
\end{chapteropening}

\section[exprimō, exprimere]{exprimō, exprimere — to press out, to express}

\label{lesson:LA-C260-exprimo}



\begin{quote}
Before the new one: say the Latin for to give up, then the Latin for to exchange.
\end{quote}

\subsection*{You'll want to know: exprimō, exprimere}
\textbf{exprimō, exprimere} — \textquotedblleft{}to press out, to express\textquotedblright{}.

\begin{cousinweb}
\textbf{exprimō} is \emph{ex}, \textquotedblleft{}out\textquotedblright{}, with \emph{premō}, \textquotedblleft{}to press\textquotedblright{} (Lewis \& Short), its e weakened to i: to press out. English \emph{express} comes from it.
\end{cousinweb}

\begin{grammarlens}[title={What you've built}]
One more word for this chapter.
\end{grammarlens}

\subsection*{Your turn}
\begin{itemize}
\raggedright
  \item \emph{Say it:} \emph{exprimō, exprimere}
  \item \emph{Say it:} \emph{exprimō, exprimere}, once more
  \item \emph{Say it:} say \emph{exprimō, exprimere}
  \item \emph{Recall:} say the Latin for to give up, then the Latin for to exchange, then say \emph{exprimō, exprimere} again
\end{itemize}

\begin{tcolorbox}[breakable,colback=teal!4,colframe=teal!35!black,title={Before you move on}]
What does exprimō, exprimere mean? (\textquotedblleft{}To press out, to express\textquotedblright{}.) Say it once more.
\end{tcolorbox}
\section[alō, alere]{alō, alere — to feed, to nourish}

\label{lesson:LA-C260-alo}



\begin{quote}
Before the new one: say the Latin for to exchange, then the Latin for to press out.
\end{quote}

\subsection*{You'll want to know: alō, alere}
\textbf{alō, alere} — \textquotedblleft{}to feed, to nourish\textquotedblright{}.

\begin{cousinweb}
\textbf{alō} is inherited, usually traced to Proto-Indo-European \emph{*h2el-}, \textquotedblleft{}to grow, to nourish\textquotedblright{} (EtymDB 2.1). Its old participle is \emph{altus}, \textquotedblleft{}grown tall, high\textquotedblright{}.
\end{cousinweb}

\begin{grammarlens}[title={What you've built}]
One more word for this chapter.
\end{grammarlens}

\subsection*{Your turn}
\begin{itemize}
\raggedright
  \item \emph{Say it:} \emph{alō, alere}
  \item \emph{Say it:} \emph{alō, alere}, once more
  \item \emph{Say it:} say \emph{alō, alere}
  \item \emph{Recall:} say the Latin for to exchange, then the Latin for to press out, then say \emph{alō, alere} again
\end{itemize}

\begin{tcolorbox}[breakable,colback=teal!4,colframe=teal!35!black,title={Before you move on}]
What does alō, alere mean? (\textquotedblleft{}To feed, to nourish\textquotedblright{}.) Say it once more.
\end{tcolorbox}
\section[compōnō, compōnere]{compōnō, compōnere — to put together, to arrange}

\label{lesson:LA-C260-compono}



\begin{quote}
Before the new one: say the Latin for to press out, then the Latin for to feed.
\end{quote}

\subsection*{You'll want to know: compōnō, compōnere}
\textbf{compōnō, compōnere} — \textquotedblleft{}to put together, to arrange\textquotedblright{}.

\begin{cousinweb}
\textbf{compōnō} is \emph{con-}, \textquotedblleft{}together\textquotedblright{}, with \emph{pōnō}, \textquotedblleft{}to put\textquotedblright{} (EtymDB 2.1). English \emph{compose} and \emph{component} come from it.
\end{cousinweb}

\begin{grammarlens}[title={What you've built}]
One more word for this chapter.
\end{grammarlens}

\subsection*{Your turn}
\begin{itemize}
\raggedright
  \item \emph{Say it:} \emph{compōnō, compōnere}
  \item \emph{Say it:} \emph{compōnō, compōnere}, once more
  \item \emph{Say it:} say \emph{compōnō, compōnere}
  \item \emph{Recall:} say the Latin for to press out, then the Latin for to feed, then say \emph{compōnō, compōnere} again
\end{itemize}

\begin{tcolorbox}[breakable,colback=teal!4,colframe=teal!35!black,title={Before you move on}]
What does compōnō, compōnere mean? (\textquotedblleft{}To put together, to arrange\textquotedblright{}.) Say it once more.
\end{tcolorbox}
\section[animadvertō, animadvertere]{animadvertō, animadvertere — to notice}

\label{lesson:LA-C260-animadverto}



\begin{quote}
Before the new one: say the Latin for to feed, then the Latin for to put together.
\end{quote}

\subsection*{You'll want to know: animadvertō, animadvertere}
\textbf{animadvertō, animadvertere} — \textquotedblleft{}to notice\textquotedblright{}.

\begin{cousinweb}
\textbf{animadvertō} is contracted from \emph{animum advertō}, \textquotedblleft{}I turn the mind towards\textquotedblright{} (Lewis \& Short): \emph{animus}, \textquotedblleft{}mind\textquotedblright{}, and \emph{advertō}, \textquotedblleft{}to turn to\textquotedblright{}.
\end{cousinweb}

\begin{grammarlens}[title={What you've built}]
One more word for this chapter.
\end{grammarlens}

\subsection*{Your turn}
\begin{itemize}
\raggedright
  \item \emph{Say it:} \emph{animadvertō, animadvertere}
  \item \emph{Say it:} \emph{animadvertō, animadvertere}, once more
  \item \emph{Say it:} say \emph{animadvertō, animadvertere}
  \item \emph{Recall:} say the Latin for to feed, then the Latin for to put together, then say \emph{animadvertō, animadvertere} again
\end{itemize}

\begin{tcolorbox}[breakable,colback=teal!4,colframe=teal!35!black,title={Before you move on}]
What does animadvertō, animadvertere mean? (\textquotedblleft{}To notice\textquotedblright{}.) Say it once more.
\end{tcolorbox}
\section[ērubēscō, ērubēscere]{ērubēscō, ērubēscere — to blush}

\label{lesson:LA-C260-erubesco}



\begin{quote}
Before the new one: say the Latin for to put together, then the Latin for to notice.
\end{quote}

\subsection*{You'll want to know: ērubēscō, ērubēscere}
\textbf{ērubēscō, ērubēscere} — \textquotedblleft{}to blush\textquotedblright{}.

\begin{cousinweb}
\textbf{ērubēscō} is \emph{ex}, \textquotedblleft{}out\textquotedblright{}, with \emph{rubēscō}, \textquotedblleft{}to grow red\textquotedblright{} (EtymDB 2.1), from \emph{ruber}, \textquotedblleft{}red\textquotedblright{}.
\end{cousinweb}

\begin{grammarlens}[title={What you've built}]
One more word for this chapter.
\end{grammarlens}

\subsection*{Your turn}
\begin{itemize}
\raggedright
  \item \emph{Say it:} \emph{ērubēscō, ērubēscere}
  \item \emph{Say it:} \emph{ērubēscō, ērubēscere}, once more
  \item \emph{Say it:} say \emph{ērubēscō, ērubēscere}
  \item \emph{Recall:} say the Latin for to put together, then the Latin for to notice, then say \emph{ērubēscō, ērubēscere} again
\end{itemize}

\begin{tcolorbox}[breakable,colback=teal!4,colframe=teal!35!black,title={Before you move on}]
What does ērubēscō, ērubēscere mean? (\textquotedblleft{}To blush\textquotedblright{}.) Say it once more.
\end{tcolorbox}
